\clearpage\begingroup
\SourceLeaf{337}\HistoricalHead{332}{L}{PRAKTISCHE UND SPHÄRISCHE ASTRONOMIE. BRIEFWECHSEL.}
\begin{center}\textsc{Bessel} an \textsc{Olbers}, 8. Februar 1821.\end{center}
\noindent\hspace*{\parindent}\{. . . . . . Ich glaube es wird Ihnen angenehm sein, das Verhalten der \textsc{Pond}schen Polardistanzen gegen die, welche mein Instrument giebt, genau kennen zu lernen. . . . . . .

. . . . . . Dass \textsc{Ponds} Polardistanzen also sämmtlich zu klein sind, darüber ist nicht der geringste Zweifel, da auch die meinigen noch um etwa \(1''\times(\cos\varphi+\sin\mathrm{ZD})\) vergrössert werden müssen. . . . . . .\}

\begin{center}\textsc{Gauss} an \textsc{Bessel}, 11. März 1821.\end{center}
\noindent\hspace*{\parindent}Ihren . . . . . . Brief vom 28. Dec[ember] würde ich schon früher beantwortet haben, wenn ich nicht gewünscht hätte, die Resultate meiner Beobachtungen der Fundamentalsterne von 1820 mit beizufügen. Erst jetzt hat \textsc{Harding} die vor einem Jahre angefangene Hülfstafel vollendet, und ich eile daher, Ihnen jetzt diese Resultate anzuzeigen, doch bloss für Sie, indem ich mir deren baldige Bekanntmachung vorbehalte. Bei der Gewichtsbestimmung ist auf die Menge der Beob[achtungen], woraus in jeder Periode der Ort des Pols bestimmt war, Rücksicht genommen und also das Gewicht Einer Zenithdistanz, als solcher, als Einheit angenommen. Die Differenzen mit den Zahlen, die Sie mir aus Ihren Beob[achtungen] angezeigt haben, sind so klein, und etwas Constantes so wenig bestimmt durchscheinend, dass unsre Instrumente beinahe als übereinstimmend angesehen werden mögen. . . . . . .

[Folgt eine Tabelle, deren Angaben unten in der Zusammenstellung Seite 346 verwendet worden sind.]

\begin{center}\textsc{Bessel} an \textsc{Olbers}, 9. April 1821.\end{center}
\noindent\hspace*{\parindent}\{. . . . . . Die constante Differenz zwischen den Declinationen südlicher Sterne, hier und in Göttingen beobachtet, scheint nun gänzlich verschwunden zu sein; wenigstens stimmen die mir von \textsc{Gauss} jetzt mitgetheilten Angaben für die Fundamentalsterne fast vollkommen und weichen nur um einige Zehntelsekunden \(+\) und \(-\) voneinander ab. Man wird dies durch Wiederholung der Beobachtungen noch weiter treiben können; allein das Hauptaugenmerk muss immer auf die constanten Fehler gerichtet sein, indem
\clearpage\SourceLeaf{338}\HistoricalHead{333}{R}{BEOBACHTUNGEN AM REICHENBACHSCHEN MERIDIANKREISE.}
\noindent die Erkenntnis derselben der Bestimmung erst Interesse und Werth geben muss. . . . . . .\}

\begin{center}\textsc{Bessel} an \textsc{Gauss}, 19. April 1821.\end{center}
\noindent\hspace*{\parindent}\{Es ist mir auffallend gewesen, aus Ihrem verehrten Briefe vom 11. März zu sehen, dass unsere Declinationen der Fundamentalsterne jetzt übereinstimmen, da dieses doch früher nicht der Fall zu sein schien. Es wird aber dadurch gezeigt, dass es immer misslich ist, auf wenige und stets in \textls[120]{einer} Jahreszeit gemachte Beob[achtungen] ein Resultat zu gründen, selbst wenn die zufälligen Beobachtungsfehler sehr klein sind. Wie unsere gemeinschaftlichen Beob[achtungen] dieser Sterne nun mit der Biegung, welche doch bei unseren Instrumenten verschieden zu sein scheint, vereinigt werden können, muss die Folge lehren. Ich theile Ihnen hier noch die Differenzen für die Sterne mit, welche ich Ihnen früher nicht angab: (Beobb. von 1820)[\begingroup\renewcommand{\thefootnote}{*)}\footnotemark\renewcommand{\thefootnote}{[*)}\footnotetext{Diese sind auf Seite 346 verwendet worden.]}\endgroup]. . . . . . .\}

\begin{center}\textsc{Olbers} an \textsc{Gauss}, 22. Januar 1824.\end{center}
\noindent\hspace*{\parindent}\{. . . . . . Bei Gelegenheit des geringen Unterschiedes, im Mittel \(1\overset{\prime\prime}{,}89\), der Münchener Polhöhe fiel mir ein, [ein]mal gehört zu haben (von wem weiss ich nicht mehr, von \textsc{Bessel} selbst war es nicht), dass \textsc{Bessel} Ihre Polhöhe von Göttingen um \(1\tfrac{1}{2}''\) zu gross halte. Wenn B[\textsc{essel}] dies wirklich glaubt oder geäussert hat, so kann dies doch wohl nur daher rühren, dass Sie aus hinreichenden Gründen \textsc{Bessels} Refraction nicht gebraucht haben. . . . .\}
\clearpage\endgroup
